\documentclass[11pt,a4paper]{article}

\usepackage[utf8]{inputenc}
\usepackage[T1]{fontenc}
\usepackage[french]{babel}
\usepackage{lmodern}
\usepackage[margin=2cm]{geometry}
\usepackage{xcolor}
\usepackage{tikz}
\usepackage{booktabs}
\usepackage{tabularx}
\usepackage{array}
\usepackage{siunitx}
\usepackage{amsmath}
\usepackage{amssymb}
\usepackage{pifont}
\usepackage{enumitem}
\usepackage{listings}
\usepackage[colorlinks=true, linkcolor=blue!50!black, urlcolor=blue!60!black]{hyperref}

\usetikzlibrary{arrows.meta, positioning, shapes.geometric, calc, fit, backgrounds}

\definecolor{chw}{HTML}{B71C1C}
\definecolor{cdsp}{HTML}{1565C0}
\definecolor{ca53}{HTML}{2E7D32}
\definecolor{cusb}{HTML}{6A1B9A}
\definecolor{cnpu}{HTML}{E65100}
\definecolor{cfx}{HTML}{00838F}
\definecolor{cbuf}{HTML}{FFB300}
\definecolor{ccomp}{HTML}{455A64}
\definecolor{codebg}{HTML}{F5F5F5}
\definecolor{cbug}{HTML}{C62828}
\definecolor{cok}{HTML}{2E7D32}

\tikzset{
  hw/.style={draw, fill=chw!15, rounded corners=2pt, font=\sffamily\footnotesize, align=center, minimum height=7mm, inner sep=3pt},
  dsp/.style={draw, fill=cdsp!15, rounded corners=2pt, font=\sffamily\footnotesize, align=center, minimum height=7mm, inner sep=3pt},
  a53/.style={draw, fill=ca53!15, rounded corners=2pt, font=\sffamily\footnotesize, align=center, minimum height=7mm, inner sep=3pt},
  pcm/.style={draw, fill=yellow!30, font=\sffamily\scriptsize, minimum height=5mm, rounded corners=1pt, align=center},
  fx/.style={draw, fill=cfx!15, rounded corners=2pt, font=\sffamily\footnotesize, align=center, minimum height=7mm, inner sep=3pt},
  usb/.style={draw, fill=cusb!15, rounded corners=2pt, font=\sffamily\footnotesize, align=center, minimum height=7mm, inner sep=3pt},
  arr/.style={-{Stealth[length=2mm]}, thick},
  arrmulti/.style={-{Stealth[length=2mm]}, line width=1.2pt}
}

\lstset{
  basicstyle=\ttfamily\footnotesize,
  backgroundcolor=\color{codebg},
  frame=single,
  rulecolor=\color{black!30},
  xleftmargin=0.5em,
  xrightmargin=0.5em,
  breaklines=true,
  columns=flexible,
  keepspaces=true,
  tabsize=2,
  inputencoding=utf8,
  extendedchars=true,
  literate={é}{{\'e}}1 {è}{{\`e}}1 {à}{{\`a}}1 {ê}{{\^e}}1 {â}{{\^a}}1 {ô}{{\^o}}1 {ù}{{\`u}}1 {û}{{\^u}}1 {î}{{\^\i}}1 {ï}{{\"\i}}1 {ç}{{\c c}}1 {É}{{\'E}}1 {È}{{\`E}}1 {À}{{\`A}}1 {→}{{$\to$}}1 {↓}{{$\downarrow$}}1 {←}{{$\leftarrow$}}1 {✓}{{\checkmark}}1 {✗}{{\texttimes}}1
}

\title{\textbf{Plateforme Mastering Live Debix}\\[6pt]
  \Large Spécification Phase 1a --- \textbf{V4.1}\\[4pt]
  \large Console numérique 8 canaux end-to-end\\
  \normalsize Intègre investigation critic V4 (6 workers, 854s)}
\author{Document de référence technique}
\date{2026-04-22 --- rev V4.1}

\begin{document}
\maketitle
\tableofcontents
\clearpage

% =======================================================================
\section{Ce qui change depuis V4}

V4 a été soumis à une nouvelle investigation (6 workers : gemini-3-pro 102s, minimax 186s, glm-5.1 395s, qwen 435s, kimi 477s, claude-code 844s). \textbf{Verdict unanime : V4 NE COMPILE PAS tel quel}. V4.1 applique les corrections convergentes (claude-code a fait l'analyse la plus précise en lisant directement les patterns upstream).

\subsection{Bugs V4 corrigés en V4.1}

\begin{center}
\begin{tabularx}{\textwidth}{llX}
\toprule
\textbf{\#} & \textbf{Sévérité} & \textbf{Correction V4.1} \\
\midrule
B1 & CRITICAL & \texttt{SOF\_DRC\_MAX\_BLOB\_SIZE} / \texttt{SOF\_MULTIBAND\_DRC\_MAX\_BLOB\_SIZE} sont des \texttt{\#define C} invisibles en m4. Remplacés par \texttt{1024} littéral (pattern upstream \texttt{pipe-drc-enabled-capture.m4:31}). \\
B2 & CRITICAL & \texttt{W\_PIPELINE(N\_PCMC/N\_PCMP, ...)} ajouté en V4 est \textbf{redondant et buggué}. Le scheduler widget est créé par \texttt{pipe-dai-capture.m4} via \texttt{DAI\_ADD} sur le même pipeline\_id. Ajouter un W\_PIPELINE côté PCM crée un 2e scheduler $\to$ comportement indéfini. \textbf{Supprimé} en V4.1. \\
B3 & CRITICAL & \texttt{multiband\_drc\_coef\_default.m4} est en mode ACTIF (emp\_deemp=1). Remplacé par \texttt{multiband\_drc\_coef\_passthrough.m4} qui existe upstream et est un vrai bypass. \\
B4 & HIGH & \texttt{drc\_coef\_default.m4} est aussi en mode actif (\texttt{enabled=1}) et \textbf{aucun blob passthrough DRC n'existe upstream}. Pire : même avec \texttt{enabled=0}, \texttt{drc\_s32\_default} applique un pre-delay de 256 frames ($\sim$5.3 ms). Conséquence : T5 sample-perfect \textbf{impossible avec DRC dans la chaîne}. V4.1 adopte 3 stratégies (voir \S\ref{sec:t5}). \\
B5 & HIGH & kcontrol PGA \texttt{"DEF\_PGA\_VOLUME"} ne matche pas le nom \texttt{SectionControlMixer} généré par \texttt{C\_CONTROLMIXER} = \texttt{"PIPELINE\_ID Master Capture Volume"}. Corrigé. \\
B6 & MEDIUM & \texttt{CONTROLMIXER} channel map : \texttt{VOLUME\_CHANNEL\_MAP} upstream (FL/FR) suffit --- en mode linked, le PGA applique un gain unique aux 8 canaux. Pas besoin de 8 KCONTROL\_CHANNEL explicites. \\
B7 & MEDIUM & \texttt{W\_PGA} periods : capture = \texttt{(2, DAI\_PERIODS, ...)}, playback = \texttt{(DAI\_PERIODS, 2, ...)} --- cohérence upstream \texttt{pipe-volume-capture.m4:47} et \texttt{pipe-volume-playback.m4}. \\
\bottomrule
\end{tabularx}
\end{center}

\subsection{Ce qui est confirmé et inchangé}

\begin{itemize}[leftmargin=*, itemsep=1pt]
  \item Architecture console 8ch end-to-end, les 4 PCMs, phasage
  \item Tous \texttt{SCHEDULE\_TIME\_DOMAIN\_DMA} via \texttt{DAI\_ADD} (pas de TIMER)
  \item Pas de cross-pipeline en Phase 1a.1 MVP
  \item Pattern \texttt{W\_VENDORTUPLES} / \texttt{W\_DATA} pour \texttt{DEF\_PGA\_CONF} validé par gemini-3-pro
  \item \texttt{CONFIG\_COMP\_MIXER=y} pour Phase 1a.3 (mixer limité 2 sources)
  \item Root cause V2.7 \texttt{-110} : \texttt{PIPELINE\_DEMUX\_3} fantôme + guard \texttt{is\_same\_sched} SCHEDULE\_TIME\_DOMAIN mismatch. Pas "W\_PIPELINE manquant" (erreur V3).
\end{itemize}

% =======================================================================
\section{Architecture cible (identique V3/V4)}

\begin{center}
\begin{tikzpicture}[node distance=4mm and 8mm, scale=0.85, transform shape]

\node[hw] (mics) {8 mics\\4× TAC5212 ADC};
\node[hw, below=14mm of mics] (hp) {HP/monitors\\4× TAC5212 DAC};

\node[dsp, right=16mm of mics, minimum width=22mm] (effin) {Effets IN\\8 canaux};
\node[pcm, right=10mm of effin] (cap) {SAI\_Capture\\ALSA 8ch cap};
\node[dsp, below=5mm of effin] (rout) {Routing 8×8};

\node[dsp, below=5mm of rout, minimum width=26mm] (mixer) {MIXER 8ch};
\node[pcm, right=10mm of mixer] (play) {SAI\_Playback\\ALSA 8ch play};
\node[usb, below=4mm of play] (ios) {iOS 2IN\\(USB gadget)};

\node[dsp, below=6mm of mixer, minimum width=22mm] (effout) {Effets OUT\\8 canaux};
\node[pcm, right=10mm of effout] (hpmon) {HP\_Monitor\\ALSA 8ch cap};
\node[usb, below=4mm of hpmon] (iosout) {iOS 2OUT\\(USB gadget)};

\draw[arrmulti] (mics.east) -- (effin.west);
\draw[arrmulti] (effin.east) -- (cap.west);
\draw[arrmulti] (effin.south) -- (rout.north);
\draw[arrmulti] (rout.south) -- (mixer.north);
\draw[arrmulti] (play.west) -- (mixer.east);
\draw[arrmulti] (ios.west) -| node[above, font=\tiny, pos=0.3]{upmix 2→8} (mixer.south east);
\draw[arrmulti] (mixer.south) -- (effout.north);
\draw[arrmulti] (effout.east) -- (hpmon.west);
\draw[arrmulti] (effout.south east) |- node[above, font=\tiny, pos=0.8]{downmix 8→2} (iosout.west);
\draw[arrmulti] (effout.west) -| (hp.north);

\end{tikzpicture}
\end{center}

% =======================================================================
\section{Phasage V4.1}

\begin{center}
\begin{tabularx}{\textwidth}{llX}
\toprule
\textbf{Phase} & \textbf{Livrables} & \textbf{Objectif} \\
\midrule
\textbf{1a.1 MVP} & 2 pipelines 8ch autonomes, effets linked (mbdrc+pga+drc), pas de cross-pipeline & Prouver passthrough + effets, trigger stable sans \texttt{-110} \\
1a.2 & Tap HP\_Monitor via \textbf{option co-scheduled cross-pipeline same sched\_comp} & Capture post-effets OUT pour NPU \\
1a.3 & Routing 8$\times$8 + mixer 2 sources (mics routés + Playback) + injection SAI\_Playback & iOS absorbé via routing slots 0-1, pas via mixer \\
1a.4 & USB gadget iOS 2IN/2OUT & iOS connectable \\
2 & USB ASIO gadget 8IN/8OUT & DAW sur host externe \\
3 & Passage linked $\to$ per-channel & 8 instances effets \\
\bottomrule
\end{tabularx}
\end{center}

% =======================================================================
\section{Phase 1a.1 MVP --- squelette m4 final V4.1}

\subsection{Architecture}

\begin{center}
\begin{tikzpicture}[node distance=4mm and 8mm]

\node[hw] (mics) {mics\\(SAI RX 8ch)};
\node[dsp, right=10mm of mics, minimum width=40mm, fill=cdsp!20] (effin) {mbdrc 8ch → pga 8ch → drc 8ch};
\node[pcm, right=10mm of effin] (cap) {SAI\_Capture\\ALSA cap 8ch};

\node[pcm, below=22mm of mics] (play) {SAI\_Playback\\ALSA play 8ch};
\node[dsp, right=10mm of play, minimum width=40mm, fill=cdsp!20] (effout) {mbdrc 8ch → pga 8ch → drc 8ch};
\node[hw, right=10mm of effout] (hp) {SAI TX 8ch\\→ HP};

\draw[arrmulti] (mics) -- (effin);
\draw[arrmulti] (effin) -- (cap);
\draw[arrmulti] (play) -- (effout);
\draw[arrmulti] (effout) -- (hp);

\node[below=1mm of effin, font=\itshape\small]{Effets IN (linked, 1 gain × 8ch)};
\node[below=1mm of effout, font=\itshape\small]{Effets OUT (linked, 1 gain × 8ch)};

\end{tikzpicture}
\end{center}

\subsection{\texttt{sof-imx8mp-tac5212-masterV41.m4} (top-level)}

\begin{lstlisting}[language={}, basicstyle=\ttfamily\scriptsize]
#
# Topology V4.1 Phase 1a.1 MVP : console 8ch end-to-end, linked effects.
# 2 pipelines 8ch autonomes, SCHEDULE_TIME_DOMAIN_DMA via DAI_ADD.
# Le scheduler widget W_PIPELINE est créé par pipe-dai-{capture,playback}.m4.
#
include(`utils.m4')
include(`dai.m4')
include(`pipeline.m4')
include(`sai.m4')
include(`pcm.m4')
include(`buffer.m4')
include(`common/tlv.m4')
include(`sof/tokens.m4')
include(`platform/imx/imx8.m4')

# PIPE 1 : SAI7 RX 8ch -> mbdrc -> pga -> drc -> SAI_Capture host
PIPELINE_PCM_ADD(sof/sof-imx8mp-tac5212-pipe-effects-capture.m4,
    1, 0, 8, s32le,
    1000, 0, 0,
    48000, 48000, 48000)

# PIPE 6 : SAI_Playback host -> mbdrc -> pga -> drc -> SAI7 TX 8ch
PIPELINE_PCM_ADD(sof/sof-imx8mp-tac5212-pipe-effects-playback.m4,
    6, 1, 8, s32le,
    1000, 0, 0,
    48000, 48000, 48000)

# DAI_ADD crée aussi le W_PIPELINE pour pipeline_id via pipe-dai-*.m4
DAI_ADD(sof/pipe-dai-capture.m4,
    1, SAI, 7, tac5212-hifi,
    PIPELINE_SINK_1, 2, s32le,
    1000, 0, 0, SCHEDULE_TIME_DOMAIN_DMA)

DAI_ADD(sof/pipe-dai-playback.m4,
    6, SAI, 7, tac5212-hifi,
    PIPELINE_SOURCE_6, 2, s32le,
    1000, 0, 0, SCHEDULE_TIME_DOMAIN_DMA)

PCM_CAPTURE_ADD(SAI_Capture, 0, PIPELINE_PCM_1)
PCM_PLAYBACK_ADD(SAI_Playback, 1, PIPELINE_PCM_6)

DAI_CONFIG(SAI, 7, 0, tac5212-hifi,
    SAI_CONFIG(DSP_A, SAI_CLOCK(mclk, 12288000, codec_mclk_in),
        SAI_CLOCK(bclk, 12288000, codec_consumer),
        SAI_CLOCK(fsync, 48000, codec_consumer),
        SAI_TDM(8, 32, 255, 255),
        SAI_CONFIG_DATA(SAI, 7, 0)))
\end{lstlisting}

\subsection{\texttt{sof-imx8mp-tac5212-pipe-effects-capture.m4}}

\begin{lstlisting}[language={}, basicstyle=\ttfamily\scriptsize]
# Capture 8ch : DAI SAI RX -> B3 -> MBDRC -> B2 -> PGA -> B1 -> DRC -> B0 -> HOST
# Pattern direct : pipe-volume-capture.m4 + pipe-drc-enabled-capture.m4 + mbdrc.

include(`utils.m4')
include(`buffer.m4')
include(`pcm.m4')
include(`dai.m4')
include(`bytecontrol.m4')
include(`mixercontrol.m4')
include(`pipeline.m4')
include(`multiband_drc.m4')
include(`pga.m4')
include(`drc.m4')

# PGA volume kcontrol (nom complet = "<PIPELINE_ID> Master Capture Volume")
C_CONTROLMIXER(Master Capture Volume, PIPELINE_ID,
    CONTROLMIXER_OPS(volsw,
        256 binds the mixer control to volume get/put handlers, 256, 256),
    CONTROLMIXER_MAX(, 32),
    false,
    CONTROLMIXER_TLV(TLV 32 steps from -64dB to 0dB for 2dB, vtlv_m64s2),
    Channel register and shift for Front Left/Right,
    VOLUME_CHANNEL_MAP)

# PGA vendor tuples (pattern pipe-volume-capture.m4)
define(DEF_PGA_TOKENS, concat(`pga_tokens_', PIPELINE_ID))
define(DEF_PGA_CONF,   concat(`pga_conf_',   PIPELINE_ID))
W_VENDORTUPLES(DEF_PGA_TOKENS, sof_volume_tokens,
LIST(`      ', `SOF_TKN_VOLUME_RAMP_STEP_TYPE  "0"'
             , `SOF_TKN_VOLUME_RAMP_STEP_MS    "250"'))
W_DATA(DEF_PGA_CONF, DEF_PGA_TOKENS)

# MBDRC byte control + blob BYPASS (multiband_drc_coef_passthrough.m4 existe upstream)
define(MULTIBAND_DRC_priv,    concat(`mbdrc_bytes_',   PIPELINE_ID))
define(MY_MULTIBAND_DRC_CTRL, concat(`mbdrc_control_', PIPELINE_ID))
include(`multiband_drc_coef_passthrough.m4')
C_CONTROLBYTES(MY_MULTIBAND_DRC_CTRL, PIPELINE_ID,
    CONTROLBYTES_OPS(bytes,
        258 binds the control to bytes get/put handlers, 258, 258),
    CONTROLBYTES_EXTOPS(
        258 binds the control to bytes get/put handlers, 258, 258),
    , , ,
    CONTROLBYTES_MAX(, 1024),
    ,
    MULTIBAND_DRC_priv)

# DRC byte control + blob (default ACTIF : pour T5 bit-perfect, voir section T5)
define(DRC_priv,    concat(`drc_bytes_',   PIPELINE_ID))
define(MY_DRC_CTRL, concat(`drc_control_', PIPELINE_ID))
include(`drc_coef_default.m4')
C_CONTROLBYTES(MY_DRC_CTRL, PIPELINE_ID,
    CONTROLBYTES_OPS(bytes,
        258 binds the control to bytes get/put handlers, 258, 258),
    CONTROLBYTES_EXTOPS(
        258 binds the control to bytes get/put handlers, 258, 258),
    , , ,
    CONTROLBYTES_MAX(, 1024),
    ,
    DRC_priv)

# Host PCM endpoint
W_PCM_CAPTURE(PCM_ID, SAI Capture, 0, 2, SCHEDULE_CORE)

# Widgets : mbdrc -> pga -> drc
W_MULTIBAND_DRC(0, PIPELINE_FORMAT, 2, 2, SCHEDULE_CORE,
    LIST(`      ', "MY_MULTIBAND_DRC_CTRL"))
W_PGA(0, PIPELINE_FORMAT, 2, DAI_PERIODS, DEF_PGA_CONF, SCHEDULE_CORE,
    LIST(`      ', "PIPELINE_ID Master Capture Volume"))
W_DRC(0, PIPELINE_FORMAT, 2, 2, SCHEDULE_CORE,
    LIST(`      ', "MY_DRC_CTRL"))

# Buffers : B0 host side, B3 DAI side
W_BUFFER(0, COMP_BUFFER_SIZE(2,
    COMP_SAMPLE_SIZE(PIPELINE_FORMAT), PIPELINE_CHANNELS,
    COMP_PERIOD_FRAMES(PCM_MAX_RATE, SCHEDULE_PERIOD)),
    PLATFORM_HOST_MEM_CAP)
W_BUFFER(1, COMP_BUFFER_SIZE(2,
    COMP_SAMPLE_SIZE(PIPELINE_FORMAT), PIPELINE_CHANNELS,
    COMP_PERIOD_FRAMES(PCM_MAX_RATE, SCHEDULE_PERIOD)),
    PLATFORM_COMP_MEM_CAP)
W_BUFFER(2, COMP_BUFFER_SIZE(2,
    COMP_SAMPLE_SIZE(PIPELINE_FORMAT), PIPELINE_CHANNELS,
    COMP_PERIOD_FRAMES(PCM_MAX_RATE, SCHEDULE_PERIOD)),
    PLATFORM_COMP_MEM_CAP)
W_BUFFER(3, COMP_BUFFER_SIZE(DAI_PERIODS,
    COMP_SAMPLE_SIZE(DAI_FORMAT), PIPELINE_CHANNELS,
    COMP_PERIOD_FRAMES(PCM_MAX_RATE, SCHEDULE_PERIOD)),
    PLATFORM_DAI_MEM_CAP)

# PAS DE W_PIPELINE ici (bug V4 B2).
# Scheduler widget fourni par pipe-dai-capture.m4 lors de DAI_ADD.

# Graph DAPM : capture = flux DAI -> HOST, dapm(sink, source)
P_GRAPH(pipe-effects-capture, PIPELINE_ID,
    LIST(`      ',
    `dapm(N_PCMC(PCM_ID), N_BUFFER(0))',
    `dapm(N_BUFFER(0), N_DRC(0))',
    `dapm(N_DRC(0), N_BUFFER(1))',
    `dapm(N_BUFFER(1), N_PGA(0))',
    `dapm(N_PGA(0), N_BUFFER(2))',
    `dapm(N_BUFFER(2), N_MULTIBAND_DRC(0))',
    `dapm(N_MULTIBAND_DRC(0), N_BUFFER(3))'))

# Exports pour DAI_ADD
indir(`define', concat(`PIPELINE_SINK_', PIPELINE_ID), N_BUFFER(3))
indir(`define', concat(`PIPELINE_PCM_',  PIPELINE_ID), SAI Capture PCM_ID)

PCM_CAPABILITIES(SAI Capture PCM_ID, CAPABILITY_FORMAT_NAME(PIPELINE_FORMAT),
    PCM_MIN_RATE, PCM_MAX_RATE, 2, PIPELINE_CHANNELS, 2, 16,
    192, 16384, 65536, 65536)

undefine(`MY_MULTIBAND_DRC_CTRL') undefine(`MULTIBAND_DRC_priv')
undefine(`MY_DRC_CTRL') undefine(`DRC_priv')
undefine(`DEF_PGA_TOKENS') undefine(`DEF_PGA_CONF')
\end{lstlisting}

\subsection{\texttt{sof-imx8mp-tac5212-pipe-effects-playback.m4}}

\begin{lstlisting}[language={}, basicstyle=\ttfamily\scriptsize]
# Playback 8ch : HOST -> B0 -> MBDRC -> B1 -> PGA -> B2 -> DRC -> B3 -> DAI SAI TX
# Pattern direct : pipe-multiband-drc-playback.m4 + pga + drc chainage.

include(`utils.m4')
include(`buffer.m4')
include(`pcm.m4')
include(`dai.m4')
include(`bytecontrol.m4')
include(`mixercontrol.m4')
include(`pipeline.m4')
include(`multiband_drc.m4')
include(`pga.m4')
include(`drc.m4')

C_CONTROLMIXER(Master Playback Volume, PIPELINE_ID,
    CONTROLMIXER_OPS(volsw,
        256 binds the mixer control to volume get/put handlers, 256, 256),
    CONTROLMIXER_MAX(, 32),
    false,
    CONTROLMIXER_TLV(TLV 32 steps from -64dB to 0dB for 2dB, vtlv_m64s2),
    Channel register and shift for Front Left/Right,
    VOLUME_CHANNEL_MAP)

define(DEF_PGA_TOKENS, concat(`pga_tokens_', PIPELINE_ID))
define(DEF_PGA_CONF,   concat(`pga_conf_',   PIPELINE_ID))
W_VENDORTUPLES(DEF_PGA_TOKENS, sof_volume_tokens,
LIST(`      ', `SOF_TKN_VOLUME_RAMP_STEP_TYPE  "2"'
             , `SOF_TKN_VOLUME_RAMP_STEP_MS    "20"'))
W_DATA(DEF_PGA_CONF, DEF_PGA_TOKENS)

# MBDRC bypass (blob passthrough upstream)
define(MULTIBAND_DRC_priv,    concat(`mbdrc_bytes_',   PIPELINE_ID))
define(MY_MULTIBAND_DRC_CTRL, concat(`mbdrc_control_', PIPELINE_ID))
include(`multiband_drc_coef_passthrough.m4')
C_CONTROLBYTES(MY_MULTIBAND_DRC_CTRL, PIPELINE_ID,
    CONTROLBYTES_OPS(bytes,
        258 binds the control to bytes get/put handlers, 258, 258),
    CONTROLBYTES_EXTOPS(
        258 binds the control to bytes get/put handlers, 258, 258),
    , , , CONTROLBYTES_MAX(, 1024), ,
    MULTIBAND_DRC_priv)

# DRC default blob (actif, voir section T5)
define(DRC_priv,    concat(`drc_bytes_',   PIPELINE_ID))
define(MY_DRC_CTRL, concat(`drc_control_', PIPELINE_ID))
include(`drc_coef_default.m4')
C_CONTROLBYTES(MY_DRC_CTRL, PIPELINE_ID,
    CONTROLBYTES_OPS(bytes,
        258 binds the control to bytes get/put handlers, 258, 258),
    CONTROLBYTES_EXTOPS(
        258 binds the control to bytes get/put handlers, 258, 258),
    , , , CONTROLBYTES_MAX(, 1024), ,
    DRC_priv)

# Host playback endpoint
W_PCM_PLAYBACK(PCM_ID, SAI Playback, 2, 0, SCHEDULE_CORE)

# Widgets : mbdrc -> pga (W_PGA periods inversés pour playback) -> drc
W_MULTIBAND_DRC(0, PIPELINE_FORMAT, 2, 2, SCHEDULE_CORE,
    LIST(`      ', "MY_MULTIBAND_DRC_CTRL"))
W_PGA(0, PIPELINE_FORMAT, DAI_PERIODS, 2, DEF_PGA_CONF, SCHEDULE_CORE,
    LIST(`      ', "PIPELINE_ID Master Playback Volume"))
W_DRC(0, PIPELINE_FORMAT, 2, 2, SCHEDULE_CORE,
    LIST(`      ', "MY_DRC_CTRL"))

# Buffers : B0 host side, B3 DAI side
W_BUFFER(0, COMP_BUFFER_SIZE(2,
    COMP_SAMPLE_SIZE(PIPELINE_FORMAT), PIPELINE_CHANNELS,
    COMP_PERIOD_FRAMES(PCM_MAX_RATE, SCHEDULE_PERIOD)),
    PLATFORM_HOST_MEM_CAP)
W_BUFFER(1, COMP_BUFFER_SIZE(2,
    COMP_SAMPLE_SIZE(PIPELINE_FORMAT), PIPELINE_CHANNELS,
    COMP_PERIOD_FRAMES(PCM_MAX_RATE, SCHEDULE_PERIOD)),
    PLATFORM_COMP_MEM_CAP)
W_BUFFER(2, COMP_BUFFER_SIZE(2,
    COMP_SAMPLE_SIZE(PIPELINE_FORMAT), PIPELINE_CHANNELS,
    COMP_PERIOD_FRAMES(PCM_MAX_RATE, SCHEDULE_PERIOD)),
    PLATFORM_COMP_MEM_CAP)
W_BUFFER(3, COMP_BUFFER_SIZE(DAI_PERIODS,
    COMP_SAMPLE_SIZE(DAI_FORMAT), PIPELINE_CHANNELS,
    COMP_PERIOD_FRAMES(PCM_MAX_RATE, SCHEDULE_PERIOD)),
    PLATFORM_DAI_MEM_CAP)

# PAS DE W_PIPELINE ici (scheduler fourni par pipe-dai-playback.m4)

# Graph DAPM : playback = flux HOST -> DAI, dapm(sink, source)
P_GRAPH(pipe-effects-playback, PIPELINE_ID,
    LIST(`      ',
    `dapm(N_BUFFER(0), N_PCMP(PCM_ID))',
    `dapm(N_MULTIBAND_DRC(0), N_BUFFER(0))',
    `dapm(N_BUFFER(1), N_MULTIBAND_DRC(0))',
    `dapm(N_PGA(0), N_BUFFER(1))',
    `dapm(N_BUFFER(2), N_PGA(0))',
    `dapm(N_DRC(0), N_BUFFER(2))',
    `dapm(N_BUFFER(3), N_DRC(0))'))

# Exports pour DAI_ADD
indir(`define', concat(`PIPELINE_SOURCE_', PIPELINE_ID), N_BUFFER(3))
indir(`define', concat(`PIPELINE_PCM_',    PIPELINE_ID), SAI Playback PCM_ID)

PCM_CAPABILITIES(SAI Playback PCM_ID, CAPABILITY_FORMAT_NAME(PIPELINE_FORMAT),
    PCM_MIN_RATE, PCM_MAX_RATE, 2, PIPELINE_CHANNELS, 2, 16,
    192, 16384, 65536, 65536)

undefine(`MY_MULTIBAND_DRC_CTRL') undefine(`MULTIBAND_DRC_priv')
undefine(`MY_DRC_CTRL') undefine(`DRC_priv')
undefine(`DEF_PGA_TOKENS') undefine(`DEF_PGA_CONF')
\end{lstlisting}

\subsection{Board config Kconfig (inchangé V4)}

\begin{lstlisting}[language={}, basicstyle=\ttfamily\scriptsize]
CONFIG_IMX8M=y
CONFIG_HAVE_AGENT=n
CONFIG_FORMAT_CONVERT_HIFI3=n
CONFIG_KPB_FORCE_COPY_TYPE_NORMAL=n

# Phase 1a.1 MVP
CONFIG_COMP_DRC=y
CONFIG_COMP_VOLUME=y
CONFIG_COMP_CROSSOVER=y
CONFIG_COMP_MULTIBAND_DRC=y
CONFIG_COMP_MUX=y

# Ajouté V4 pour Phase 1a.3
CONFIG_COMP_MIXER=y
\end{lstlisting}

% =======================================================================
\section{Test T5 bit-perfect --- limite découverte}\label{sec:t5}

L'investigation V4 (claude-code) a relevé un point important : même avec \texttt{enabled=0} côté DRC, la fonction \texttt{drc\_s32\_default} (\texttt{drc\_generic.c:744-750}) applique un \textbf{pre-delay de 256 frames} ($\sim$5.3 ms @48k). Le null test sample-perfect \textbf{n'est donc pas atteignable avec DRC dans la chaîne}.

Trois stratégies T5 :

\begin{enumerate}[leftmargin=*, itemsep=1pt]
  \item \textbf{T5a compensé} : loopback HW, null test avec compensation de latence ($\sim$256 frames DRC + buffers). Tolérance $\pm$1 LSB. Faisable, pragmatique.
  \item \textbf{T5b sans DRC} : topology variante \texttt{masterV41-noDRC} qui skip le widget DRC dans la chaîne. Null test sample-perfect. Utilisé uniquement pour la gate V4.1.
  \item \textbf{T5c blob patché} : créer localement un \texttt{drc\_coef\_zerolatency.m4} avec \texttt{enabled=0} ET \texttt{pre\_delay\_time=0}. Forcerait \texttt{drc\_default\_pass = audio\_stream\_copy()}. Pratique long terme mais patch custom.
\end{enumerate}

\textbf{Choix V4.1} : T5a en premier passage (pragma\-tique), T5c comme option long terme si T5a peine à converger.

% =======================================================================
\section{Tests Phase 1a.1 MVP V4.1}

\begin{enumerate}[leftmargin=*, itemsep=2pt]
  \item \textbf{T0 Compile} : \texttt{m4 -I. -I m4 -I sof -I common masterV41.m4 | alsatplg -c - -o masterV41.tplg} sans erreur, aucun widget orphelin. \textbf{Gate obligatoire} avant deploy.
  \item \textbf{T0a Widget count} : \texttt{alsatplg --dump masterV41.tplg} $\to$ doit afficher 2 W\_PIPELINE (créés par pipe-dai-*), pas 4 (preuve que V4.1 n'ajoute pas de scheduler redondant). Total $\sim$22 widgets : 2 PCMs + 8 buffers + 6 effets + 2 DAI + 2 W\_PIPELINE + 2 schedulers scheduling\_plat.
  \item \textbf{T1 Load} : kernel charge sans \texttt{-22} ni \texttt{-110}. Carte \texttt{softac5212tdm} + 2 PCMs.
  \item \textbf{T2 Playback} : \texttt{aplay -D plughw:softac5212tdm,1 -c 8 -f S32\_LE tone.wav}.
  \item \textbf{T3 Capture} : \texttt{arecord -D plughw:softac5212tdm,0 -c 8 -f S32\_LE -d 3 mics.wav}.
  \item \textbf{T4 Full-duplex} : T2 + T3 simultanés 60 s, pas de xrun.
  \item \textbf{T5 Bit-perfect} : stratégie T5a (loopback HW + compensation latence).
  \item \textbf{T6 Contrôles} : \texttt{amixer cset name='1 Master Capture Volume' 32} change le gain. Les 2 bytes controls MBDRC/DRC visibles.
  \item \textbf{T7 Charge CPU} : lecture via sof-logger ou debugfs \texttt{/sys/kernel/debug/sof/etrace} des cycles \texttt{kcps} ou \texttt{cpc}.
\end{enumerate}

% =======================================================================
\section{Phase 1a.2 --- stratégie tap HP\_Monitor}

L'investigation V4 a confirmé : \textbf{aucun exemple upstream} de \texttt{W\_PCM\_CAPTURE} comme sink d'un widget dans une pipeline PLAYBACK. Le framework IPC3 ne l'interdit pas explicitement (\texttt{host\_legacy.c} détermine direction par widget type, pas par position dans la chaîne), mais le pattern est non testé.

\textbf{Stratégie V4.1 retenue} : \textbf{option co-scheduled cross-pipeline same sched\_comp}.

\begin{itemize}[leftmargin=*, itemsep=1pt]
  \item Pipeline playback (PIPE 6) reste avec effets OUT. Un \texttt{W\_MUXDEMUX} inséré en fin alimente 2 sinks : B3 (vers DAI TX) et B4 (vers pipeline capture HP\_Monitor).
  \item Pipeline HP\_Monitor (PIPE 7) : pipeline capture \textbf{piggyback scheduling} sur PIPE 6 via \texttt{PIPELINE\_SCHED\_COMP\_6}. Le \texttt{sched\_comp} partagé passe le guard \texttt{is\_same\_sched} (\texttt{pipeline-stream.c:505}).
  \item DAPM cross-pipeline : \texttt{SectionGraph dapm(PIPELINE\_SINK\_7, N\_MUXDEMUX\_via\_PIPELINE\_DEMUX\_6)} --- en s'assurant que \texttt{PIPELINE\_DEMUX\_6} est bien exporté par \texttt{pipe-effects-playback-with-tap.m4} (bug V2.7 à ne pas reproduire).
\end{itemize}

Fallback : loopback HW (câble TAC OUT $\to$ TAC IN) si le pattern cross-pipeline tap reste instable.

% =======================================================================
\section{Questions pour l'investigation V4.1}

\begin{enumerate}[leftmargin=*, itemsep=2pt]
  \item \textbf{Gate compile V4.1} : le squelette \texttt{pipe-effects-capture.m4} (section 4.3) et \texttt{pipe-effects-playback.m4} (section 4.4) compilent-ils \textbf{ligne par ligne} via \texttt{m4 | alsatplg} ? Vérifier notamment :
    \begin{itemize}[leftmargin=*]
      \item \texttt{multiband\_drc\_coef\_passthrough.m4} existe bien dans \texttt{sof/tools/topology/topology1/m4/} et produit un blob valide ?
      \item \texttt{C\_CONTROLMIXER} args 6 et 7 ("Channel register..." + \texttt{VOLUME\_CHANNEL\_MAP}) : l'argument 6 (\texttt{Channel register and shift for Front Left/Right}) est-il bien une \textbf{string libre non-quotée} comme dans le pattern upstream, ou faut-il le quoter/échapper ?
      \item Le nom kcontrol référencé \texttt{"PIPELINE\_ID Master Capture Volume"} -- est-ce que m4 expand bien \texttt{PIPELINE\_ID} ici au moment de l'inclusion du fichier par \texttt{PIPELINE\_PCM\_ADD}, ou faut-il une autre syntaxe ?
    \end{itemize}
  \item \textbf{Pattern V\_PGA periods capture/playback} : capture = \texttt{(2, DAI\_PERIODS)}, playback = \texttt{(DAI\_PERIODS, 2)}. C'est cohérent avec \texttt{pipe-volume-capture.m4:47} et \texttt{pipe-volume-playback.m4:33}. Mais est-ce que W\_MULTIBAND\_DRC et W\_DRC en 8ch doivent aussi utiliser \texttt{DAI\_PERIODS} du côté DAI, ou \texttt{2, 2} suffit ?
  \item \textbf{Widget W\_PIPELINE absent --- confirmation} : le pattern upstream \texttt{pipe-drc-enabled-capture.m4}, \texttt{pipe-volume-capture.m4}, \texttt{pipe-multiband-drc-playback.m4} n'a \textbf{pas} de W\_PIPELINE. Celui-ci est créé dans \texttt{pipe-dai-capture.m4:XX} / \texttt{pipe-dai-playback.m4:YY} via \texttt{DAI\_ADD}. Confirmer en lisant ces deux fichiers que le W\_PIPELINE créé utilise bien le même \texttt{PIPELINE\_ID} que \texttt{PIPELINE\_PCM\_ADD} $\to$ un seul scheduler widget total.
  \item \textbf{Co-scheduling Phase 1a.2 tap} : pour que le tap HP\_Monitor fonctionne via \texttt{PIPELINE\_SCHED\_COMP\_6} piggyback, quelle est la syntaxe exacte ? 13e arg de \texttt{PIPELINE\_PCM\_ADD} ? Y a-t-il des exemples upstream (regarder \texttt{sof-smart-amplifier.m4} ou \texttt{echo-ref}) ? Est-ce que la pipeline capture co-scheduled doit elle-même avoir \texttt{SCHEDULE\_TIME\_DOMAIN\_DMA} ou hériter via le piggyback ?
  \item \textbf{DRC pre-delay dans T5} : confirmer la valeur exacte de \texttt{CONFIG\_DRC\_MAX\_PRE\_DELAY\_FRAMES} (256 ? autre?) dans notre fork et que \texttt{drc\_s32\_default} applique \textbf{toujours} ce pre-delay même avec \texttt{enabled=0}. Existe-t-il une variante \texttt{drc\_func} qui soit "enabled=0 + no pre-delay" ? Sinon, confirmer que la seule solution bit-perfect est de \textbf{retirer le widget DRC du graph} (stratégie T5b).
  \item \textbf{MBDRC passthrough blob} : \texttt{multiband\_drc\_coef\_passthrough.m4} existe-t-il dans notre branche \texttt{feature/sdma-ap2ap-phase1} ? Si oui, fournir le contenu (num\_bands, emp\_deemp, crossover coeffs). Si non, quel fork amont l'a (\texttt{sof main}, NXP, Intel) et faut-il le copier ?
  \item \textbf{Red flags résiduels} : après toutes les corrections V4.1, reste-t-il des pièges non anticipés pour le premier build ? Notamment autour de :
    \begin{itemize}[leftmargin=*]
      \item \texttt{DAI\_PERIODS} défini dynamiquement par \texttt{DAI\_ADD} --- le widget défini dans le fichier pipeline y a-t-il accès au moment de l'expansion m4 ?
      \item \texttt{SCHEDULE\_PERIOD} = 1000 us $\to$ 48 frames period $\to$ \texttt{COMP\_PERIOD\_FRAMES(PCM\_MAX\_RATE, SCHEDULE\_PERIOD)} = 48, conforme aux buffer sizes mbdrc \texttt{buf\_src[PLATFORM\_MAX\_CHANNELS]} ?
      \item Le nom du topology file \texttt{masterV41.m4} doit être ajouté à \texttt{CMakeLists.txt} pour que la compilation west/bitbake le prenne.
    \end{itemize}
\end{enumerate}

% =======================================================================
\section{Livrables Phase 1a.1 MVP V4.1}

\begin{itemize}[leftmargin=*, itemsep=1pt]
  \item \texttt{sof/tools/topology/topology1/sof-imx8mp-tac5212-masterV41.m4}
  \item \texttt{sof/tools/topology/topology1/sof/sof-imx8mp-tac5212-pipe-effects-capture.m4}
  \item \texttt{sof/tools/topology/topology1/sof/sof-imx8mp-tac5212-pipe-effects-playback.m4}
  \item Patch \texttt{sof/app/boards/imx8mp\_evk\_mimx8ml8\_adsp.conf} : \texttt{+CONFIG\_COMP\_MIXER=y}
  \item Patch \texttt{sof/tools/topology/topology1/CMakeLists.txt} : ajouter \texttt{masterV41.m4} à la liste
  \item \texttt{meta-local/recipes-kernel/linux/files/test-phase1a1-mvp-v41.sh} (tests T0--T7)
  \item Suppression des fichiers V2.7 buggés : \texttt{masterlite.m4}, \texttt{pipe-masterchain.m4}, \texttt{pipe-mastertap.m4}, \texttt{pipe-sai-merge.m4}, \texttt{mux\_route\_sai\_merge.m4}, \texttt{demux\_route\_default.m4}
\end{itemize}

% =======================================================================
\section{Protocole}

\begin{itemize}[leftmargin=*, itemsep=1pt]
  \item Aucune modification source SOF sans \texttt{critic\_analyze} préalable (mode DOUX, fenêtre 2000 s post-analyze).
  \item Un changement = un commit = validation utilisateur.
  \item Gate \texttt{m4 | alsatplg} OBLIGATOIRE avant tout deploy sur board.
  \item Si T0 échoue, retour à \texttt{critic\_analyze} avec l'erreur exacte \textbf{avant} de modifier le squelette.
  \item Backup kernel+firmware baseline avant premier deploy (rescue).
\end{itemize}

\end{document}
